% Historical source (removed by author request): references/_OLD_Übungsaufgaben/Übungsblatt 2 - Lineare Gleichungssysteme/Blatt 2.tex:18
\begin{ExerciseBlock}
% Compact display gaps keep the seven systems with the first-page sheet head.
\setlength{\abovedisplayskip}{6pt}
\setlength{\belowdisplayskip}{6pt}
\setlength{\abovedisplayshortskip}{3pt}
\setlength{\belowdisplayshortskip}{3pt}
\ExerciseHeading{01}{Gauß-Verfahren}{linear_systems_gaussian_elimination}
Lösen Sie die folgenden Gleichungssysteme mit Hilfe des Gauß-Verfahrens. Entscheiden Sie, ob die Gleichungssysteme eine eindeutige, keine oder unendlich viele Lösungen haben.
\begin{enumerate}[label=\alph*),itemsep=0pt,parsep=0pt,topsep=0pt]
	\item 
	\begin{equation*}
	\begin{split}
		2x + 4y  &= 2 \\
 		8x + 4y &= 8
	\end{split}
	\end{equation*}

	\item 
	\begin{equation*}
	\begin{split}
		2x + 4y  &= 3 \\
 		8x + 16y &= 12
	\end{split}
	\end{equation*}

	\item 
	\begin{equation*}
	\begin{split}
		2x + 4y  &= 3 \\
 		8x + 16y &= 6
	\end{split}
	\end{equation*}

	\item 
	\begin{equation*}
	\begin{split}
		x + y + z &= 42 \\
 		2x+y-z &= 21 \\
 		5x+2y+3z &= -21
	\end{split}
	\end{equation*}

	\item 
	\begin{equation*}
	\begin{split}
		x_1 + 2 x_2 + 3 x_3 + x_4 &= 0 \\
 		x_1 + 3 x_2 + 2 x_3 + x_4 &= 1 \\
 		2x_1 +  x_2 - x_3 + 4 x_4 &= 1 \\
 		x_1 + 6 x_2 + 4 x_3 + 2 x_4 &= 0
	\end{split}
	\end{equation*}

	\item 
	\begin{equation*}
	\begin{split}
		x_1 + 2 x_2 &= 0 \\
 		2 x_1 + 3 x_2 - 2 x_3 &= 0 \\
 		x_1 + x_2 - 2 x_3 &= 1 \\
	\end{split}
	\end{equation*}

	\item 
	\begin{equation*}
	\begin{split}
		x_1 + 2 x_2 &= 0 \\
 		2 x_1 + 3 x_2 - 2 x_3 &= 0 \\
 		x_1 + x_2 - 2 x_3 &= 0 \\
	\end{split}
	\end{equation*}

\end{enumerate}

\end{ExerciseBlock}
\UseExerciseSolution{linear_systems_gaussian_elimination}
